% !TeX TS-program = xelatex
\documentclass{resume-photo}

% 字体（Overleaf 推荐用 Noto 系列，避免 SimSun 缺失）
\usepackage{xeCJK}
\setCJKmainfont{Noto Serif CJK SC}
\usepackage{fontspec}

% 紧凑排版
\usepackage{enumitem}
\setlist[itemize]{itemsep=2pt, topsep=2pt}
\newcommand{\ExpItem}[3]{%
  \noindent\textbf{#1｜#2} \hfill {\footnotesize #3}\\[-2pt]
}

\ResumeName{邓祎明}
\ResumePhoto{dengyiming.jpg}

\begin{document}

\ResumeContacts{
  177-7918-0948,%
  \ResumeUrl{mailto:dengyiming.ken@foxmail.com}{dengyiming@std.uestc.edu.cn},%
  \textnormal{电子科技大学 | 生物医学工程 · 硕士 | 2002-09}%
}

\ResumeTitle


% ====================== 简介 ======================
\noindent 电子科技大学生物医学工程硕士在读，具备计算机视觉算法研发、端侧模型落地与跨团队工程协作经验。曾在旷视科技手机业务部门实习，独立推进注视检测模型从研发到交付，在复杂场景准确率与部署效果上取得显著提升。当前研究方向为跨设备眼动跟踪与无障碍人机交互，聚焦跨域泛化、实时交互与系统级优化，具备从算法设计到产品化实现的完整实践能力。


% ====================== 教育 ======================
\section{教育背景}
\ResumeItem
{电子科技大学}
[\textnormal{生命科学与技术学院 | 生物医学工程 | 硕士}]
[2024.09—2027.06(预计)]

\ResumeItem
{桂林电子科技大学}
[\textnormal{生命与环境科学学院 | 生物医学工程 | 学士}]
[2020.09—2024.06]

% ====================== 技能 ======================
\section{专业技能}
\begin{itemize}
  \item 熟悉 Python、PyTorch、MegEngine 等深度学习开发框架，具备目标检测与目标跟踪项目经验。
  \item 熟悉 YOLOv8、ByteTrack、BoT-SORT 等视觉算法，能够完成检测/跟踪任务建模、训练与评估。
  \item 具备端侧模型优化与部署经验，熟悉 QAT/PTQ 量化流程及跨硬件平台一致性问题分析。
  \item 具备数据工程与模型工程实践能力，能够搭建数据与模型管理流程并支持多项目并行迭代。
  \item 熟悉实时眼动跟踪与人机交互系统开发，具备跨设备校准、误差优化与交互设计能力。
\end{itemize}

% ====================== 科研 ======================
\section{科研经历}

\ResumeItem{跨设备重定向增强方法研究}
[]
[2024.09—至今]

\begin{itemize}
  \item \textbf{问题背景}: 跨设备视线估计存在尺寸差异引起的分布偏移问题，导致注视点在不同视野设备间映射不稳定。
  \item \textbf{研究内容}: 1) 提出基于重定向数据增强的跨设备泛化方法，将小视野注视数据扩增到大视野设备；2) 结合个性化注视点校准与多点快速采样策略，提升系统跨用户与跨设备鲁棒性；3) 构建跨设备评测流程，对比不同重定向策略在泛化精度与稳定性上的表现。
  \item \textbf{相关成果}: 目前已形成相关论文一篇，处于 TIP 投稿阶段。
\end{itemize}


% ====================== 实习 ======================
\section{实习经历}

\ResumeItem{旷视科技｜移动业务事业部}
[\textnormal{算法实习生}]
[2024.04—2024.10]

\begin{itemize}
  \item 独立完成蔚来手机注视检测模型的研发与交付，负责算法研发与客户数据对接；基于 MegEngine 框架完成约 4 个月迭代优化，针对客户反馈 badcase 从数据与算法两侧改进。
  \item 通用场景准确率由 84\% 提升至 95\%，复杂场景准确率由 62\% 提升至 85\%，模型成功上线至蔚来最新款手机。
  \item 主导注视检测基座模型更新，使用全向人脸检测替代旧版四向检测，结合人脸 landmark 与仿射变换提升侧脸/歪脸场景鲁棒性；整体准确率提升约 2\%，半脸场景提升约 30\%。
  \item 完成算法基础设施建设与多项目支持，基于 nori2 搭建统一云端数据与模型管理流程，实现数据、模型、测试文件分离，训练效率提升约 150\%；并参与 vivo、荣耀等项目算法迭代。
  \item 参与跨算法、开发、测试、量化、研究院多部门协作，推进 QAT 与 PTQ 方案调试，解决量化与推理结果对齐问题，提升模型跨硬件平台一致性与可移植性。
\end{itemize}

\ResumeItem{欧洲蚊虫控制协会（EMCA）/ Entostudio}
[\textnormal{国际实习生}]
[2023.04—2023.10]

\begin{itemize}
  \item 开发基于 YOLOv8 的目标检测模型，用于识别监测视频中的蟑螂与蚂蚁个体。
  \item 开发基于 ByteTrack、BoT-SORT 的目标跟踪模型，用于捕捉视频中所有个体的运动轨迹。
  \item 设计并实现基于 Qt 的软件界面，集成检测与跟踪流程，支持标定区域内目标计数与计时。
\end{itemize}



% ====================== 项目 ====================== 
\section{项目经历}

\ResumeItem{实时眼动跟踪交互平台开发}
[]
[2024.09—至今]

\begin{itemize}
  \item \textbf{项目背景}: 为提升肢体受限用户的计算机使用能力，开发基于普通摄像头的实时眼动跟踪交互平台。
  \item \textbf{核心痛点}: 纯视觉方案在实时性、注视点精度和个体差异适配方面存在明显挑战，影响可用性与稳定性。
  \item \textbf{技术方案}: 1) 构建基于笔记本摄像头的眼动跟踪与交互闭环，实现眼动驱动的鼠标控制、键盘输入及常用操作；2) 设计个性化校准与多点快速采样流程，降低跨用户误差；3) 设计多种辅助工具与交互界面，优化操作路径与反馈机制。
  \item \textbf{实验验证}: 在标准测试场景下进行时延与精度评测，系统帧率达到 30FPS 以上，平均注视点预测误差小于 2cm。
  \item \textbf{最终效果}: 平台具备稳定的人机交互能力，可支持基础输入与常用系统操作，显著提升无障碍使用体验。
\end{itemize}


% ====================== 荣誉 ======================
\section{个人荣誉}
\begin{itemize}
  \item 以第一作者发表计算机视觉领域 SCI 论文 1 篇
  \item 申请发明专利 2 项、软件著作权 1 项
  \item 第十届全国大学生生物医学工程创新设计竞赛 三等奖
  \item 第十八届“挑战杯”全国大学生课外学术科技作品竞赛 三等奖
  \item 第九届中国国际“互联网+”大学生创新创业大赛广西赛区 金奖
\end{itemize}

\end{document}
